\section{Introduction}
\label{sec:introduction}

\subsection{Motivation}
Modern enterprise networks emit log streams from a heterogeneous mix of
systems: web servers, audit daemons, virtual private network endpoints,
name-resolution services, and authentication subsystems. Security
operations teams must reason across these streams to detect attacker
activity, attribute events to a kill-chain phase, and reconstruct a
timeline after the fact. In practice, the work is hampered because
each log type lives in its own format, storage location, and tool. A
relational database offers a single substrate where structured queries
can join evidence across sources and ground-truth labels.

This project takes a published, labeled, multi-host security log corpus
\cite{landauer2022maintainable, ait2024dataset} and asks whether a
single PostgreSQL database in third normal form can support a
meaningful cross-log analysis of a real multi-stage attack. The dataset
ships with attacker ground truth: every log line that participates in
the simulated attack chain is labeled with one or more attack-phase
identifiers, allowing us to quantify both detection coverage and the
cost of false positives.

\subsection{Problem and Research Gap}
The AIT-LDSv2.0 corpus is large and heterogeneous: 295{,}243 log lines
from eight sources on the \texttt{russellmitchell} testbed, with
61{,}862 labeled attack events across seven kill-chain phases. Most
published exploration of the dataset treats each log file
independently, or flattens the entire corpus into a single denormalized
table, losing the relational structure that makes cross-source queries
cheap. We close that gap by:
\begin{itemize}
  \item Designing a 3NF schema that explicitly separates the host
    inventory, audit event, and attack-label domains while preserving
    cross-domain provenance and foreign keys.
  \item Implementing a reproducible ETL from raw files through staging
    tables into the final 3NF schema, with Alembic-managed migrations.
  \item Demonstrating that the resulting database supports analytical
    SQL expressive enough to reconstruct attack sequences and to
    compare detection-rule effectiveness across log sources.
\end{itemize}

\subsection{Scope and Story Arc}
Across the eight log sources in the \texttt{russellmitchell} slice, the
attack chain traces a coherent sequence: a web foothold on
\texttt{intranet\_server}'s \texttt{access.log.2} during the
\texttt{initial\_access} phase, a \texttt{su} then
\texttt{sudo cat /etc/shadow} chain captured by the host's
\texttt{audit.log} during the \texttt{privilege\_escalation} phase,
and a DNS-tunneling \texttt{put} service on \texttt{internal\_share}
during the \texttt{exfiltration} phase roughly nine hours later. The
narrative through this report follows the lead-in plus climax pair on
\texttt{intranet\_server}; the exfiltration phase is reported as
``what happened next'' in Section~\ref{sec:conclusion} and reproduced
in the appendix.

\subsection{Project Goals}
The project has four explicit deliverables, each of which is reflected
in a section of this report:
\begin{enumerate}
  \item A normalized relational schema with documented PK/FK and
    normalization decisions (Section~\ref{sec:schema}).
  \item A reproducible PostgreSQL load pipeline with row-count
    verification (Section~\ref{sec:dbsetup}).
  \item A portfolio of 23 SQL query items (7 basic and 16 advanced,
    spread across 6 basic and 13 advanced \texttt{.sql} files) that
    answer concrete attack-investigation questions, including
    indexing/\texttt{EXPLAIN} performance work and a saved view
    (Section~\ref{sec:sql}).
  \item An interactive dashboard with nine visualizations that
    surfaces the query results to a non-SQL audience
    (Section~\ref{sec:dashboard}).
\end{enumerate}
